\documentclass[11pt,a4paper]{amsart}

\usepackage[T1]{fontenc}
\usepackage[utf8]{inputenc}
\usepackage{amsmath,amssymb,amsthm}
\usepackage{tikz}
\usepackage[hidelinks]{hyperref}
\usepackage{booktabs}

\theoremstyle{plain}
\newtheorem{theorem}{Theorem}[section]
\newtheorem{lemma}[theorem]{Lemma}
\newtheorem{proposition}[theorem]{Proposition}
\newtheorem{corollary}[theorem]{Corollary}
\theoremstyle{definition}
\newtheorem{definition}[theorem]{Definition}
\newtheorem{convention}[theorem]{Convention}
\newtheorem{remark}[theorem]{Remark}


\begin{document}

\title[Three conjectures on alternating plane graphs]
      {Three conjectures on alternating plane graphs}

\author{Hanyu Yang}
\email{hanyu.yang.92@gmail.com}
\urladdr{\href{https://orcid.org/0009-0005-0419-4070}{ORCID 0009-0005-0419-4070}}

\date{\today}

\subjclass[2020]{05C10, 05C30, 68R10}
\keywords{alternating plane graph, plane graph, rotation system,
computer-assisted proof}

\begin{abstract}
Alth\"ofer, Haugland, Scherer, Schneider and Van Cleemput introduced
\emph{alternating plane graphs} and closed their paper with four open problems.
We settle three of them.

Conjecture~10.1, that no alternating plane graph has exactly two distinct
vertex degrees or exactly two distinct face sizes, is proved. The argument is a
single per-edge inequality applied twice, and it is unconditional: we show that
such a graph would have to be bridgeless, so the proof does not depend on how
the definition is read at a bridge.

Conjecture~10.2, that $(3,4,5)$-alternating plane graphs exist on $n$ vertices
for every $n \ge 20$, is settled. We exhibit verified certificates for all
$26$ orders that were open, and prove a periodic capping lemma: one certificate
per residue class generates a $(3,4,5)$-alternating plane graph at order $48$
and at every order $n \ge 50$.

Conjecture~10.3, that a $3$-connected alternating plane graph exists on $n$
vertices for every $n \ge 19$, is settled. We also refute the natural shortcut
by exhibiting a $(3,4,5)$-alternating plane graph on $46$ vertices with a
separating pair.

The remaining problem --- the asymptotic distribution of the degree ratio
$v_4/v_3$ on $[1,1.5]$ --- is untouched, and we explain why the present methods
do not reach it.

Every graph asserted here is supplied as an explicit rotation system and is
accepted by three independently written verifiers; the accompanying artifact
runs $1253$ checks and requires only a Python interpreter.
\end{abstract}

\maketitle

\section{Introduction}

Alth\"ofer, Haugland, Scherer, Schneider and Van Cleemput \cite{AHSSV2015}
introduced alternating plane graphs: plane graphs in which adjacent vertices
have different degrees and adjacent faces have different sizes. The condition is
severe. It forces a global rigidity that makes even existence questions hard,
and the paper closes (\cite{AHSSV2015}, \S10, p.~362) with four open problems.
Three are settled here.

\subsection*{Results}

\begin{itemize}
  \item Conjecture~10.1 is proved, unconditionally (Theorem~\ref{thm:101}).
  \item Conjecture~10.2 is settled, by certificates together with a proved
        construction (Theorem~\ref{thm:pumping}, Corollary~\ref{cor:102}).
  \item Conjecture~10.3 is settled (Theorem~\ref{thm:103}).
  \item The asymptotic degree distribution is left open (\S\ref{sec:density}).
\end{itemize}

\subsection*{Disclosure of AI use}

This work was produced with substantial assistance from large language models:
the software, the search programs that found the certificates, the arguments,
and the exposition. The step that removes the (C2) hypothesis in
\S\ref{sec:c2}, and so makes Theorem~\ref{thm:101} unconditional, was supplied
by an AI reviewer and subsequently re-derived and verified by the author. AI
review was also used adversarially, and found errors that had been missed,
including three claims stated in earlier drafts that were false and were
withdrawn. No AI system is an author of this paper, and responsibility for every
claim in it rests with the named author. The accompanying artifact states the
extent in full.

\subsection*{On the use of computation}

Two of the three results assert the existence of specific graphs, and those
graphs are large: the certificates here run to $110$ vertices. We state them as
rotation systems, which determine a plane map up to homeomorphism, and we derive
every claimed property --- degrees, faces, face sizes, bridges, connectivity ---
rather than storing it. Nothing in the artifact records a claim that a graph
\emph{is} an alternating plane graph; three separately written decision
procedures recompute that on every run.

The proofs in \S\ref{sec:101} and \S\ref{sec:pumping} are ordinary mathematics
and can be read without a computer. Computation enters in two places: verifying
the certificates, and discharging two finite hypotheses of the capping lemma,
both stated explicitly in \S\ref{sec:pumping} so that a reader can see precisely
what is being taken from a machine.

\section{Preliminaries}\label{sec:prelim}

We follow \cite{AHSSV2015} throughout.

\begin{definition}[\cite{AHSSV2015}, Definition~2.1]\label{def:apg}
An \emph{alternating plane graph} is a simple connected plane graph in which
\begin{enumerate}
  \item[(a)] every vertex has degree at least $3$;
  \item[(b)] every face, the unbounded one included, has size at least $3$;
  \item[(c)] adjacent vertices have different degrees;
  \item[(d)] adjacent faces have different sizes.
\end{enumerate}
\end{definition}

\begin{definition}[\cite{AHSSV2015}, Definitions~3.1 and~3.4]
A $(3,4,5)$-alternating plane graph is one in which every vertex degree and
every face size lies in $\{3,4,5\}$. An \emph{$X,Y$-alternating plane graph} is
one with exactly $X$ distinct vertex degrees and exactly $Y$ distinct face
sizes.
\end{definition}

\begin{convention}[face size]\label{conv:c1}
The \emph{size} of a face is the length of its boundary walk counted with
multiplicity, that is, the number of edge-side incidences. It is not the number
of distinct vertices on the face. The two differ exactly when the walk repeats a
vertex, which happens if and only if the graph is not $2$-connected. In
particular a bridge contributes $2$ to the size of the single face carrying it.
This is the reading used in \cite{AHSSV2015} at (3.1) and (9.2), both of which
assert $\sum_s s\,f_s = 2e$.
\end{convention}

Condition (d) of Definition~\ref{def:apg} does not mention bridges, and at a
bridge one face lies on both sides of an edge, so the condition is either
violated there or vacuous there. Two readings are therefore available in
principle, but \cite{AHSSV2015} does choose between them, as we record next.

\begin{convention}[bridges]\label{conv:c2}
Say that a plane graph satisfies \emph{(C2)} if no edge has the same face on
both sides, equivalently if it has no bridge. \cite{AHSSV2015} takes (d) to
include a face adjacent to itself, so that (C2) follows from
Definition~\ref{def:apg}; this is stated explicitly on p.~339, where an
alternating plane graph is said to be ``always at least $2$-edge-connected,
since plane graph with edge connectivity $1$ contains a face that is adjacent to
itself''. Under the \emph{weak reading}, (d) constrains only pairs of distinct
faces and bridges are permitted.
\end{convention}

The proof of Conjecture~10.1 in \S\ref{sec:101} uses (C2), and we prove in
\S\ref{sec:c2} that a putative counterexample must be bridgeless under
\emph{either} reading. The result is therefore independent of the choice. We
stress that this is a robustness statement and not the repair of a gap: the
authors of \cite{AHSSV2015} state the strict reading, and we simply decline to
rely on it.

We shall use the following, which is due to \cite{AHSSV2015}.

\begin{theorem}[\cite{AHSSV2015}, Theorem~3.2]\label{thm:32}
In a $(3,4,5)$-alternating plane graph, $v_3 = f_3$, $v_4 = f_4$ and
$v_5 = f_5$, where $v_i$ and $f_i$ count vertices of degree $i$ and faces of
size $i$. Consequently $v_5 = v_3 - 4$, $E = 2V - 2$ and $F = V$.
\end{theorem}

\subsection{Rotation systems}

A plane map is encoded by its \emph{rotation system}: for each vertex, its
neighbours in clockwise order. Writing $\alpha$ for the involution exchanging
the two darts of an edge and $\sigma$ for the rotation, the faces are the orbits
of $\varphi = \sigma^{-1}\alpha$, and the size of a face is the length of its
orbit. An edge is a bridge if and only if its two darts lie in the same orbit.
Every graph in this paper is presented this way.

\section{The per-edge identity}\label{sec:identity}

The following double count is the engine of \S\ref{sec:101}. It is the
convention behind (9.2) of \cite{AHSSV2015}, isolated.

\begin{lemma}\label{lem:identity}
Let $G$ be a plane graph with $v$ vertices, $e$ edges and $f$ faces. For an edge
$a$ with ends $u,w$ and face occurrences $F^-(a)$, $F^+(a)$ on its two sides,
put
\[
  c(a) \;=\; \frac{1}{\deg u} + \frac{1}{\deg w}
           + \frac{1}{|F^-(a)|} + \frac{1}{|F^+(a)|}.
\]
Then $\sum_{a} c(a) = v + f$. If $G$ is connected, Euler's formula gives
\[
  \sum_{a} c(a) \;=\; e + 2 ,
  \qquad\text{equivalently}\qquad
  \sum_{a} \bigl(c(a) - 1\bigr) \;=\; 2 .
\]
\end{lemma}

\begin{proof}
A vertex of degree $d$ is an end of exactly $d$ edges and contributes $1/d$ each
time, hence $1$ in total; summing gives $v$. By Convention~\ref{conv:c1} a face
of size $s$ occurs exactly $s$ times as a side of an edge and contributes $1/s$
each time, hence $1$ in total; summing gives $f$.
\end{proof}

\begin{remark}
Neither count requires facial walks to be simple. Repeated vertices, cut
vertices and bridges are all permitted; a bridge simply has
$F^-(a) = F^+(a)$ and contributes $2/|F|$ from its face side. This is what makes
Lemma~\ref{lem:identity} usable in \S\ref{sec:c2}, where bridges are the object
of study.
\end{remark}

We write $x(a) = c(a) - 1$ for the \emph{excess} of an edge, so that
Lemma~\ref{lem:identity} reads $\sum_a x(a) = 2$: in a connected plane graph
some edge must have positive excess.

\section{Conjecture 10.1}\label{sec:101}

\begin{theorem}\label{thm:101}
There is no $2,Y$-alternating plane graph and no $X,2$-alternating plane graph.
\end{theorem}

We prove the two halves separately, then remove the dependence on
Convention~\ref{conv:c2} in \S\ref{sec:c2}. Both halves are the same inequality
with the roles of vertices and faces exchanged; no dual graph is constructed, so
neither invokes the $3$-edge-connectivity hypothesis of
(\cite{AHSSV2015}, Lemma~2.2).

\subsection{No $2,Y$-alternating plane graph}

Suppose $G$ is one, with degrees taking exactly the two values $d_1 < d_2$.

Adjacent vertices differ in degree and only two degrees occur, so every edge
joins a $d_1$-vertex to a $d_2$-vertex. In particular $G$ is bipartite, and
every edge contributes
\[
  \frac{1}{d_1} + \frac{1}{d_2} \;\le\; \frac13 + \frac14 \;=\; \frac{7}{12}
\]
to the vertex side, using $d_1 \ge 3$ from Definition~\ref{def:apg}(a).

Since $G$ is bipartite, every closed walk has even length, so by
Convention~\ref{conv:c1} every face has even size, hence size at least $4$. By
(C2) the two sides of an edge are distinct faces, and by
Definition~\ref{def:apg}(d) their sizes differ; being distinct even numbers at
least $4$, one is at least $4$ and the other at least $6$. So the face side
contributes at most $\tfrac14 + \tfrac16 = \tfrac{5}{12}$.

Hence $c(a) \le \tfrac{7}{12} + \tfrac{5}{12} = 1$ for every edge, so
$\sum_a x(a) \le 0$, contradicting Lemma~\ref{lem:identity}. \qed

\begin{remark}
This is (9.5) of \cite{AHSSV2015} with $1/2$ replaced by $1/d_1 \le 1/3$. The
derivation of (9.5) never used the degree $2$; \S9.1 of \cite{AHSSV2015} runs
the chain on the weak class with degrees $2$ and $k$, while \S3.2 attacks the
strong class with different machinery yielding only $V \ge 25$ and $V \ge 56$
and never bounding $f$. The two sections were not connected.
\end{remark}

\subsection{No $X,2$-alternating plane graph}\label{sec:halftwo}

Suppose $G$ is one, with face sizes taking exactly the two values $s_1 < s_2$.

By (C2) the two sides of every edge are distinct faces, and by
Definition~\ref{def:apg}(d) they have different sizes; with only two sizes
available, every edge separates an $s_1$-face from an $s_2$-face and contributes
at most $\tfrac13 + \tfrac14 = \tfrac{7}{12}$ to the face side.

We claim every degree is even. Fix a vertex $u$ of degree $d$, let its edges in
rotation order be $a_1,\dots,a_d$, and let $F_i$ be the face occupying the
corner between $a_i$ and $a_{i+1}$, indices modulo $d$. The two sides of $a_i$
are the occurrences $F_{i-1}$ and $F_i$, distinct faces by (C2), so
$|F_{i-1}| \ne |F_i|$. With only two sizes the cyclic word
$|F_1|,\dots,|F_d|$ alternates, and a cyclically alternating binary word has
even length. Hence $d$ is even.

Stating the argument over corners rather than over ``the incident faces'' is
what makes it correct at a cut vertex, where one face occupies several corners:
those corners are non-consecutive, so alternation is undisturbed.

Degrees are at least $3$ and even, hence at least $4$; adjacent vertices have
distinct even degrees, so the vertex side contributes at most
$\tfrac14+\tfrac16 = \tfrac{5}{12}$. Again $c(a) \le 1$ for every edge and
Lemma~\ref{lem:identity} is contradicted. \qed

\begin{remark}
Euler is used as $v - e + f = 2$. Disconnectedness only strengthens the
contradiction: with $c$ components $v + f = e + 1 + c$.
\end{remark}

\subsection{Calibration}

The chain can be run on the class \cite{AHSSV2015} settles, the weak $2,k$
class, where it must exclude exactly $k \ge 12$, reproducing
(\cite{AHSSV2015}, (9.10)), and must \emph{not} exclude $k \le 10$, for which
\cite{AHSSV2015} exhibits weak $2,k$-alternating plane graphs (Table~5,
Figures~8--9). It does both. Any step that over-counted would have failed this
test. The case $k = 11$ is left open by the chain, exactly as in
\cite{AHSSV2015}, which needs its Lemma~9.4 for it.

\section{Removing the dependence on (C2)}\label{sec:c2}

Everything in this section is under the weak reading of
Convention~\ref{conv:c2}: condition (d) of Definition~\ref{def:apg} constrains
only pairs of distinct faces, and bridges are permitted. We show that no
$X,2$-alternating plane graph exists even so. Since such a graph would in
particular be bridgeless, (C2) is a consequence of Definition~\ref{def:apg}
rather than a reading of it, and Theorem~\ref{thm:101} holds under both
readings. (The corresponding statement for $2,Y$ follows from a leaf-block
argument, sketched at the end of the section.)

\begin{lemma}[parity]\label{lem:parity}
Let $G$ be a plane graph whose faces admit a colouring by two colours that is
proper on every pair of distinct adjacent faces. Then for every vertex $v$,
\[
  \deg v \;-\; \bigl|\{\text{bridges at } v\}\bigr| \quad\text{is even.}
\]
\end{lemma}

\begin{proof}
Walk the darts at $v$ in rotation order and let $F_1,\dots,F_d$ be the corner
faces as in \S\ref{sec:halftwo}, so that $F_{i-1}$ and $F_i$ are the two sides
of the edge $a_i$. These two occurrences are of the same face precisely when
$a_i$ is a bridge. Hence the corner colour changes across each non-bridge edge
at $v$ and is unchanged across each bridge at $v$. The walk closes up, so the
number of changes is even, and that number is $\deg v$ minus the number of
bridges at $v$.
\end{proof}

Lemma~\ref{lem:parity} is the classical fact that a plane graph is face
$2$-colourable if and only if it is Eulerian, restated so that bridges are
permitted rather than assumed away. In an $X,2$-alternating plane graph the face
size supplies the colouring, so the lemma applies.

\begin{corollary}\label{cor:evendeg}
In an $X,2$-alternating plane graph under the weak reading, a vertex lying on no
bridge has even degree, hence degree at least $4$.
\end{corollary}

\begin{theorem}\label{thm:nobridge}
Under the weak reading there is no $X,2$-alternating plane graph. In particular
no such graph has a bridge.
\end{theorem}

\begin{proof}
Let $G$ be one, with face sizes $s_1 < s_2$, and let $B$ be its number of
bridges. Suppose first $B \ge 1$.

\smallskip
\emph{Step 0: $s_2 \ge 8$.} Let $a = uw$ be a bridge and let $A$ be the
component of $G - a$ containing $u$. Then $\deg_A u = \deg u - 1 \ge 2$, so $A$
has at least three vertices. A simple connected plane graph on at least three
vertices has every facial walk of length at least $3$: a walk of length $2$
requires parallel edges and one of length $1$ a loop. The face $F$ carrying $a$
is traced as the outer walk of $A$, then $a$, then the outer walk of the other
component, then $a$ again, so by Convention~\ref{conv:c1}
\[
  |F| \;=\; |\partial A| + |\partial B| + 2 \;\ge\; 3 + 3 + 2 \;=\; 8 .
\]
As $|F| \in \{s_1,s_2\}$ and $s_1 < s_2$, we get $s_2 \ge 8$.

\smallskip
\emph{Step 1: only an edge at a vertex of degree $3$ can have positive excess.}
Recall $s_1 \ge 3$ and, by Step 0, $s_2 \ge 8$; degrees at the two ends of an
edge are distinct by Definition~\ref{def:apg}(c). Three cases:
\begin{itemize}
  \item $a$ a bridge: both sides carry $F$, so
        $c(a) \le \tfrac13+\tfrac14+\tfrac28 = \tfrac56$ and
        $x(a) \le -\tfrac16$;
  \item $a$ not a bridge, both ends of degree at least $4$:
        $c(a) \le \tfrac14+\tfrac15+\tfrac13+\tfrac18 = \tfrac{109}{120}$ and
        $x(a) \le -\tfrac{11}{120}$;
  \item $a$ not a bridge with an end of degree $3$:
        $c(a) \le \tfrac13+\tfrac14+\tfrac13+\tfrac18 = \tfrac{25}{24}$ and
        $x(a) \le \tfrac1{24}$.
\end{itemize}
Only the third class can be positive. Note in particular that a bridge endpoint
of degree $5$ cannot make an edge positive.

\smallskip
\emph{Step 2: each bridge carries at most two positive edges.} First, a positive
edge is a $3$--$4$ edge. Its face side is at most $\tfrac13+\tfrac18 =
\tfrac{11}{24}$, so its vertex side must exceed $\tfrac{13}{24}$; among pairs of
distinct degrees at least $3$, only $\tfrac13+\tfrac14 = \tfrac{14}{24}$ does,
since already $\tfrac13+\tfrac15 = \tfrac{12.8}{24}$ falls short.

Let $a$ be an edge of positive excess and let $u$ be its end of degree $3$,
which is unique since adjacent vertices have different degrees. By
Lemma~\ref{lem:parity} the
number of bridges at $u$ has the same parity as $\deg u = 3$, so it is $1$ or
$3$; it is not $3$, since $a$ itself is a non-bridge edge at $u$. So $u$ lies on
exactly one bridge, hence on exactly two non-bridge edges, and we may send $a$
to that bridge.

The qualification matters: a vertex of degree $3$ need not carry two non-bridge
edges, since one with three bridges carries none. What is proved, and what is
needed, is that a vertex of degree $3$ \emph{incident with a non-bridge edge}
lies on exactly one bridge and two non-bridge edges.

Now let $b$ be a bridge. Its two ends have different degrees by
Definition~\ref{def:apg}(c), so at most one of them has degree $3$. Hence at
most two positive edges are sent to $b$, and writing $P$ for the number of
positive edges, $P \le 2B$.

\smallskip
\emph{Step 3.} Every non-positive edge has $x(a) \le 0$, so by
Lemma~\ref{lem:identity},
\[
  2 \;=\; \sum_a x(a) \;\le\; \frac{P}{24} - \frac{B}{6}
     \;\le\; \frac{2B}{24} - \frac{B}{6} \;=\; -\frac{B}{12} \;<\; 0 ,
\]
a contradiction.

\smallskip
Finally suppose $B = 0$. By Corollary~\ref{cor:evendeg} every degree is even and
at least $4$, so the vertex side of every edge is at most $\tfrac14+\tfrac16 =
\tfrac5{12}$; the face side is at most $\tfrac13+\tfrac14 = \tfrac7{12}$; hence
$x(a) \le 0$ for every edge and $\sum_a x(a) = 2$ is again contradicted.
\end{proof}

\begin{remark}[how much room the argument has]
The tether of Step 2 is the substance: treating the positive edges as
unconstrained gives only a lower bound on $B$ and closes nothing. It does not
have to be sharp. Ignoring the degree condition on a bridge's two ends gives
only $P \le 4B$, and even then the bound of Step 3 is $0$, which still
contradicts $2$; the sharp tether merely supplies the margin $-B/12$. That
margin also means Step 0 is not needed at full strength: with $s_2 \ge 7$ the
per-bridge total is $-\tfrac1{84}$, still negative. Some lower bound on the
bridge face is nevertheless indispensable --- at $s_2 \ge 6$ the total is
$+\tfrac1{12}$ --- and at $P \le 5B$ the argument fails outright.
\end{remark}

\begin{proposition}\label{prop:twoY-weak}
Under the weak reading there is no $2,Y$-alternating plane graph either.
\end{proposition}

\begin{proof}
Step 0 of Theorem~\ref{thm:nobridge} uses only minimum degree $3$ and
simplicity, so it applies here: a bridge's face has size at least $8$. As in
\S\ref{sec:101} the graph is bipartite, since two degrees and
Definition~\ref{def:apg}(c) force every edge to join a $d_1$-vertex to a
$d_2$-vertex; hence every closed walk has even length, a facial walk traversing
a bridge twice included, and every face has even size at least $4$. Then a
bridge has
\[
  c(a) \;\le\; \tfrac13 + \tfrac14 + \tfrac28 \;=\; \tfrac56 ,
\]
and a non-bridge edge, whose two sides are distinct faces of distinct even sizes
at least $4$ and $6$, has
$c(a) \le \tfrac13+\tfrac14+\tfrac14+\tfrac16 = 1$. So $x(a) \le 0$ for every
edge, contradicting $\sum_a x(a) = 2$.
\end{proof}

\section{The periodic capping lemma}\label{sec:pumping}

Conjecture~10.2 asks for a $(3,4,5)$-alternating plane graph on $n$ vertices for
every $n \ge 20$. Twenty-six orders were open. We supply certificates for all of
them (\S\ref{sec:artifact}) and, more usefully, prove that a single certificate
generates an infinite family, so that the result does not rest on a finite list.

\subsection{The cover}

The certificates are capped unrollings of a periodic quotient of the torus. The
relevant structure is a bi-infinite \emph{strip}, built from identical
\emph{copies}, each contributing three vertices and six edges, together with two
\emph{caps} that close it off at the ends. Writing $\mathrm{TARGET}_n$ for a
certificate of order $n$, its alignment against the cover exhibits a maximal run
of copies whose rotations are translates of one another; we call this run the
\emph{deep block} $D$.

\begin{theorem}[periodic capping lemma]\label{thm:pumping}
Fix an order $n$ whose alignment against the cover has a deep block $D$ of at
least five copies, and let $\mathrm{cut}$ be a copy at distance at least two
from each end of $D$. Then for every integer $d \ge -(|D| - 1)$ the rotation
system $\mathrm{splice}(n,d)$, obtained from $\mathrm{TARGET}_n$ by inserting
$d$ further translates of the copy $\mathrm{cut}$ (or deleting $-d$ of them when
$d < 0$), is a $(3,4,5)$-alternating plane graph on $n + 3d$ vertices.
\end{theorem}

\begin{proof}
Write $S = \mathrm{splice}(n,d)$. Every vertex of $S$ receives its rotation in
one of four ways: cap vertices verbatim, with strip neighbours above the cut
shifted by $d$; strip copies at or below the cut verbatim; strip copies above
$\mathrm{cut}+d$ verbatim from copy $c-d$; and the $d$ fresh copies from the cut
copy's rotation, translated.

The argument is local. A facial walk is traced by
$\varphi = \sigma^{-1}\alpha$, which reads only the edge partner and the
rotation predecessor at each vertex it visits. A face of a
$(3,4,5)$-alternating plane graph has at most five darts and the cover's largest
edge offset is two, so a facial walk cannot span more than three consecutive
copies. The fresh copies lie in $(\mathrm{cut}, \mathrm{cut}+d]$, and
$\mathrm{cut}$ is at distance at least two from each end of $D$, so every
five-copy window around a fresh copy is a translate of the window
$[\mathrm{cut}-2,\mathrm{cut}+2]$ of $\mathrm{TARGET}_n$; away from the fresh
block $S$ agrees with $\mathrm{TARGET}_n$ verbatim. Hence \emph{every face of
$S$ is a translated face of $\mathrm{TARGET}_n$}.

Each condition of Definition~\ref{def:apg} is now inherited, because each is
local: face sizes and vertex degrees lie in $\{3,4,5\}$; adjacent vertices have
unequal degrees and adjacent faces unequal sizes; facial walks are simple; and a
loop or parallel edge would have to sit inside a two-copy window, where $S$ and
$\mathrm{TARGET}_n$ agree. Connectivity is inherited because each copy meets the
next and every cap vertex reaches the strip as before.

Finally $S$ is spherical. A copy contributes three vertices, six edges and ---
since every face is a translated face --- three faces, and $3-6+3=0$, so
$V-E+F$ is unchanged from $\mathrm{TARGET}_n$, where it is $2$.
\end{proof}

\begin{remark}[what is machine-checked]\label{rem:hypotheses}
Two facts in the proof are not bookkeeping and are discharged by computation on
a single certificate: that no face spans more than three consecutive copies, and
that the deep block really is periodic, i.e.\ its rotations are translates of one
another. Both are finite checks. Everything else is the local argument above.
\end{remark}

\subsection{What the family reaches}

The floor $-(|D|-1)$ is where deletion would consume a copy that is not deep,
that is, where the two cap collars meet; each collar is two copies wide.
Splicing changes the order by $3d$, so a certificate reaches one residue class
modulo $3$. The floors are:

\begin{center}
\begin{tabular}{ccl}
\toprule
$n \bmod 3$ & floor & family \\
\midrule
$0$ & $48$ & $48, 51, 54, 57, \dots$ \\
$1$ & $52$ & $52, 55, 58, 61, \dots$ \\
$2$ & $50$ & $50, 53, 56, 59, \dots$ \\
\bottomrule
\end{tabular}
\end{center}

\begin{corollary}\label{cor:102}
Conjecture~10.2 holds: there is a $(3,4,5)$-alternating plane graph on $n$
vertices for every $n \ge 20$.
\end{corollary}

\begin{proof}
Orders below $46$ are the published witnesses of \cite{AHSSV2015}, decoded from the author-hosted
corpus~\cite{APGcorpus}, together with
the Section-8 closures of \S\ref{sec:103}. Orders $46$ through $56$, $67$
through $74$, $88$ through $92$, and $109$, $110$ carry the certificates of this
paper. Every remaining order is reached by Theorem~\ref{thm:pumping} from a
certificate in its residue class modulo $3$, with the floors above.
\end{proof}

\begin{remark}[the deletion direction]\label{rem:deletion}
Theorem~\ref{thm:pumping} is stated for every $d \ge -(|D|-1)$, and the floors
in the table are reached by deletion rather than insertion. The locality
argument is written above for insertion, where the fresh copies are interior to
$D$; for $d < 0$ it reads the same way with the four cases collapsing, save at
the floor itself, where a three-copy window meets both cap collars and ``every
window is a translated window'' is no longer literally true. The floor orders
are therefore also carried by their own verified certificates, and orders $57$,
$58$, $59$, $61$ and $63$ --- which the hypothesis $|D| \ge 5$ does not reach by
insertion alone --- rest on machine-verified splices rather than on the theorem
as stated. Tightening the hypothesis to cover deletion uniformly is left open;
nothing else in this paper depends on it.
\end{remark}

This is an independent construction of (\cite{AHSSV2015}, Theorem~8.1), which
gives $n \ge 111$, and of twenty-three of the twenty-six open orders, from one
certificate per residue class. Orders $46$, $47$ and $49$ lie outside the family
and rest on their certificates alone.

Consequently the certificates and Theorem~\ref{thm:pumping} together settle
\emph{every order $n \ge 46$} without reference to \cite{AHSSV2015}, subject to
Remark~\ref{rem:deletion} at the five orders noted there; Theorem~8.1 is not
needed. Orders $20$ through $45$ are inherited from that paper and are not
re-established here.

\begin{remark}[why this is not twenty-six more examples]
The splice reproduces certificates it did not start from.
$\mathrm{TARGET}_{110}$ shortened by twenty periods is $\mathrm{TARGET}_{50}$ up
to isomorphism of rotation systems, and so is $\mathrm{TARGET}_{92}$ shortened
by fourteen; $\mathrm{TARGET}_{90}$ shortened by fourteen and
$\mathrm{TARGET}_{72}$ shortened by eight both give $\mathrm{TARGET}_{48}$. Two
independently found certificates decompose into the same pair of caps and the
same period, which is what shows the cap extraction to be canonical rather than
an artefact of one alignment.
\end{remark}

\section{Conjecture 10.3}\label{sec:103}

\begin{theorem}\label{thm:103}
For every $n \ge 19$ there is a $3$-connected alternating plane graph on $n$
vertices.
\end{theorem}

The witnesses come from six sources, and no order depends on more than one of
them being correct.

\begin{center}
\begin{tabular}{ll}
\toprule
orders & source \\
\midrule
$19$ & all five alternating plane graphs on $19$ vertices \\
$17, 20, 26$--$36, 42$ & published witnesses, decoded and re-verified \\
$21, 22, 23, 25$ & published search seeds \\
$21$--$24$, $39$--$45$ & Section-8 closures, rebuilt here \\
$37, 38$ & disk surgery \\
$46$--$56$, $67$--$74$, $88$--$92$, $109$, $110$ & the certificates of \S\ref{sec:pumping} \\
$48$ and every $n \ge 50$ & the spliced family \\
\bottomrule
\end{tabular}
\end{center}

Three orders deserve comment.

\emph{Order $19$} is the one order the settled Conjecture~10.2 cannot reach:
\cite{AHSSV2015} reports no $(3,4,5)$-alternating plane graph on $18$ or $19$
vertices, so a witness must be a general alternating plane graph in the sense of
Definition~\ref{def:apg}. There are exactly five on $19$ vertices~\cite{HoG}, and all
five are $3$-connected. A single one would have sufficed.

\emph{Orders $37$ and $38$} are out of reach of the Section-8 arithmetic
$18a+19b+20c+21d+3$, because $34$ and $35$ are not sums of $18,19,20,21$; this
is why the coverage list of \cite{AHSSV2015} skips them. They are obtained by
disk surgery: excise the star of a vertex, an edge or a face from a stored
$(3,4,5)$-alternating plane graph, let the boundary degrees float, and refill
the disk.

\emph{The infinite tail} is not an artefact of a search horizon. The spliced
family is $3$-connected at \emph{every} order it produces, by the same locality
that carries Theorem~\ref{thm:pumping}; finite sweeps only check that theorem's
hypotheses.

\subsection{A shortcut that does not exist}

If every $(3,4,5)$-alternating plane graph were $3$-connected, then
Corollary~\ref{cor:102} would give Theorem~\ref{thm:103} for all $n \ge 20$ in
one step.

\begin{proposition}\label{prop:not3conn}
There is a $(3,4,5)$-alternating plane graph on $46$ vertices with a separating
pair. In particular $(3,4,5)$-alternating plane graphs need not be
$3$-connected.
\end{proposition}

The graph is supplied as a rotation system in the artifact; both verifiers
accept it, its unique separating pair is exhibited, and removing that pair
leaves two components of $22$ vertices each. Nothing in
Theorem~\ref{thm:103} depends on the proposition --- every witness is checked
individually rather than inferred from a class-wide claim --- but it is why the
connectivity of the spliced family has to be proved rather than assumed.

\section{The asymptotic degree distribution}\label{sec:density}

The fourth problem of (\cite{AHSSV2015}, \S10, p.~362) is not a counting
question but a question about the shape of a typical graph. Writing $r = v_3$,
Theorem~\ref{thm:32} confines $v_4$ to $[r-5, \tfrac32 r - 5]$, so the ratio
$v_4/v_3$ lies asymptotically in $[1, 1.5]$. The authors ask whether there is a
density function on that interval giving the asymptotic fractions of
$(3,4,5)$-alternating plane graphs for large $n$, and if so what it looks like.

We have nothing to offer on it. What we can add is a caution, from the only
data this paper produces.

All twenty-six certificates satisfy $v_5 = v_3 - 4$, as they must, and all lie
inside the interval. But they lie at its very bottom: $v_4/v_3$ ranges over
$[1.0000, 1.1250]$ across the twenty-six, and drifts \emph{downwards} with $n$,
attaining exactly $1$ at orders $50$, $56$, $74$, $92$ and $110$.

That is a selection effect, not a measurement. The certificates come from one
construction --- a capped unrolling of a periodic strip --- and
Theorem~\ref{thm:pumping} extends it by inserting copies of a single period, so
every member of the family inherits that period's degree profile. A construction
emitting one graph per order reports on the construction, not on the
distribution over all graphs of that order. The clustering at $v_4/v_3 = 1$ is
precisely what a periodic family looks like and precisely what a sample would
not.

The methods of this paper cannot reach the question, for a structural reason.
Everything proved above is an existence statement: one graph per order, or a
class shown empty. A density function requires a measure over all graphs of
order $n$, and nothing here enumerates. The class is moreover rigid in a way
that obstructs sampling: a sweep over every two-edge switch of all twenty-six
certificates finds no switch leaving a valid $(3,4,5)$-alternating plane graph,
so the one-switch neighbourhood of each is empty and local moves cannot walk the
space.

Three routes seem plausible. An exhaustive census at small orders, with
\texttt{plantri}~\cite{plantri} or an equivalent generator restricted to the degree and
face-size profile, would give the first honest histogram of $v_4/v_3$; it is a
means to a conjecture rather than to a theorem, and it is where we would start.
A structure theorem giving \emph{every} $(3,4,5)$-alternating plane graph a
cap--strip--cap decomposition with a bounded interface alphabet would make the
count a regular language and the density a transfer-matrix computation; the
obstacle is the word \emph{every}, since we have a construction and not a
classification. Failing either, a bijection with an already-counted family of
plane maps would serve, and the identities of Theorem~\ref{thm:32} are the
natural hook.

\section{The artifact}\label{sec:artifact}

Every graph asserted in this paper is supplied as an explicit rotation system in
the accompanying artifact, in a documented JSON format with a dependency-free
reference reader. No file records a claim that a graph \emph{is} an alternating
plane graph: degrees, faces, face sizes, bridges, connectivity and both
alternation conditions are recomputed on every run by three separately written
decision procedures. Certificates can also be exported to \texttt{planar\_code}
for checking by \texttt{plantri}~\cite{plantri} or House of Graphs~\cite{HoG}.

The artifact runs $1253$ checks and needs only a Python interpreter; the
mathematics uses no third-party library. Two classes of check are load-bearing:
that each certificate is accepted by every verifier, and that each source in
Theorem~\ref{thm:103} is necessary, in the sense that removing it reopens
orders. Each gate is paired with a control that must fail, so that a check
cannot pass by being vacuous.

The two computations that the proofs genuinely rest on are the finite hypotheses
of Remark~\ref{rem:hypotheses}.

\begin{thebibliography}{9}

\bibitem{AHSSV2015}
I.~Alth\"ofer, J.~K. Haugland, K.~Scherer, F.~Schneider and N.~Van Cleemput,
\emph{Alternating plane graphs},
Ars Math. Contemp. \textbf{8} (2015), 337--363.
\href{https://doi.org/10.26493/1855-3974.584.09a}{doi:10.26493/1855-3974.584.09a}.

\bibitem{plantri}
G.~Brinkmann and B.~D. McKay,
\emph{Fast generation of planar graphs} (expanded edition),
MATCH Commun. Math. Comput. Chem. \textbf{58} (2007), 323--357.

\bibitem{HoG}
K.~Coolsaet, S.~D'hondt and J.~Goedgebeur,
\emph{House of Graphs 2.0: A database of interesting graphs and more},
Discrete Appl. Math. \textbf{325} (2023), 97--107.
\href{https://doi.org/10.1016/j.dam.2022.10.013}{doi:10.1016/j.dam.2022.10.013}.

\bibitem{APGcorpus}
I.~Alth\"ofer, \emph{Alternating plane graphs} --- status page and
machine-readable corpus of published witnesses,
\url{https://althofer.de/alternating-plane-graphs.html} and
\url{https://althofer.de/apg/apgs/}, accessed 1 September 2026.

\end{thebibliography}

\end{document}
